\documentclass[11pt]{article}

\usepackage[margin=1in]{geometry}
\usepackage{amsmath, amssymb, amsthm}
\usepackage{enumitem}

\newtheorem{lemma}{Lemma}
\newtheorem{theorem}{Theorem}

\title{Bounding the OLS Slope of a Time Series\\
       When You Only Have Bucket Summaries}
\author{}
\date{}

\begin{document}
\maketitle

\section{The problem, concretely}

We have a time series of raw points $(x_{ij}, y_{ij})$. We want one
number: the ordinary least-squares (OLS) slope of a contiguous stretch
of that series, the segment. Written out, that slope is

\[
\beta \;=\; \frac{S_{xy}}{S_{xx}},
\qquad
S_{xx}=\sum_{i,j}(x_{ij}-\bar X)^2,
\qquad
S_{xy}=\sum_{i,j}(x_{ij}-\bar X)(y_{ij}-\bar Y),
\]

where $\bar X,\bar Y$ are the means of all the points in the segment.

\paragraph{What we actually have.}
We do \emph{not} have the raw points. The store has already chopped the
segment into $n$ consecutive \emph{buckets} and thrown the raw points
away, keeping for each bucket $i$ only a five-number summary:

\begin{itemize}[noitemsep]
  \item its time span $[a_i,b_i]$ (exact calendar boundaries),
  \item how many points it holds, $c_i$,
  \item the smallest and largest $y$ in it, $\ell_i$ and $h_i$,
  \item the exact mean of its $y$ values, $\mu_i$.
\end{itemize}

That is everything. We never see an individual $x_{ij}$ or $y_{ij}$
again.

\paragraph{What we want to produce.}
Because the raw points are gone, we cannot compute $\beta$ exactly --
many different point clouds are consistent with the same bucket
summaries, and they have different slopes. So instead of a single
number we will produce an \textbf{interval} $[\beta^-,\beta^+]$ that is
\emph{guaranteed} to contain the true slope, whatever the (unknown) raw
points were, as long as they are consistent with the summaries. The
whole proof is the construction of that interval and the argument that
it really does trap $\beta$.

\paragraph{The one source of uncertainty.}
For a single bucket, the summary pins down a lot but not everything.
The $y$ side is well controlled: every $y_{ij}$ lies in $[\ell_i,h_i]$,
and we know the bucket's exact $y$-mean $\mu_i$. The $x$ side is the
problem. We know every $x_{ij}$ lies in the time span $[a_i,b_i]$, but
we have \emph{no} information about where inside that span the points
fall. They could be bunched at the left edge, bunched at the right,
spread out -- the summary does not say. This single unknown -- the
horizontal placement of points inside each bucket -- is the entire
source of slack in everything that follows. Keep an eye on it.

\section{Setting up the bookkeeping}

Fix notation per bucket $i$:

\[
m_i=\frac{a_i+b_i}{2}\ \ (\text{the midpoint}),
\qquad
w_i=b_i-a_i\ \ (\text{the width}),
\qquad
N=\sum_i c_i\ \ (\text{total points}).
\]

The true (but unknown) average $x$-position of the points in bucket
$i$ is

\[
\bar x_i \;=\; \frac{1}{c_i}\sum_j x_{ij}.
\]

We do not know $\bar x_i$. All the summary tells us is that it lies
somewhere in the bucket's span $[a_i,b_i]$.

\paragraph{Why the midpoint --- a minimax choice.}
To compute anything we must pick a single representative $x$-position
for bucket $i$ to stand in for the unknown $\bar x_i$. Call our choice
$p_i\in[a_i,b_i]$. Whatever we pick, the error we will carry is
$\bar x_i-p_i$, and an adversary -- choosing the worst point cloud
consistent with the summary -- may place the true $\bar x_i$ anywhere
in $[a_i,b_i]$. So the worst-case size of our error is

\[
\max_{\bar x_i\in[a_i,b_i]}\bigl|\bar x_i-p_i\bigr|
=\max\bigl(p_i-a_i,\ b_i-p_i\bigr).
\]

We get to pick $p_i$ \emph{before} the adversary picks $\bar x_i$, so
we should pick the $p_i$ that makes this worst case as small as
possible. The maximum of the two distances $p_i-a_i$ and $b_i-p_i$ is
smallest exactly when they are equal,

\[
p_i-a_i=b_i-p_i
\quad\Longleftrightarrow\quad
p_i=\frac{a_i+b_i}{2}=m_i,
\]

i.e.\ the \textbf{midpoint}. Any other point in the bucket makes one of
the two distances longer, and the adversary simply places $\bar x_i$ at
that far end, giving a strictly larger guaranteed error. The midpoint
is therefore not a convenience -- it is the unique choice that
minimizes the worst-case $x$-error, and with it that worst case is
exactly half the bucket width:

\[
\delta_i \;:=\; \bar x_i - m_i,
\qquad
|\delta_i|\le \frac{w_i}{2}.
\tag{$\star$}
\]

Inequality $(\star)$ is the \emph{only} thing we know about the
$x$-placement, and the whole bound is built by pushing $(\star)$
through the algebra. Because the midpoint achieves the smallest
possible constant ($w_i/2$) on the right of $(\star)$, every error term
that later inherits this constant -- $E_P$, $E_\downarrow$,
$E_\uparrow$ -- is as tight as the available information allows. Pick
any other representative and all of them grow.

The $y$ side needs no such choice: the summary stores the
\emph{exact} bucket mean $\mu_i$, so the bucket's $y$-centroid is
exactly $\mu_i$ (write $\bar y_i=\mu_i$) -- an equality, not a bound.

\paragraph{Quantities we can actually compute.}
We cannot compute $\bar X$ (it needs the unknown $\bar x_i$). But with
the midpoint as each bucket's $x$-representative, and the exact mean as
each bucket's $y$-representative, we can compute the count-weighted
surrogates

\[
\hat X=\frac1N\sum_i c_i\,m_i,
\qquad
\hat Y=\frac1N\sum_i c_i\,\mu_i.
\]

These are our stand-ins for the unknown true means $\bar X,\bar Y$.
(A second, technical reason the midpoint is convenient: it makes the
sums $\sum_i c_i(\mu_i-\hat Y)$ and $\sum_i c_i u_i$ collapse to zero,
which is what kills the first-order error terms in Lemmas~B and~C.
That cancellation actually works for any affine choice of
representative; the minimax argument above is the reason the midpoint
in particular is the \emph{right} affine choice.) The first job is to
measure how wrong the stand-ins are.

\section{Step 1 --- the means (Lemma A)}

\begin{lemma}\label{lem:means}
$\bar Y=\hat Y$ exactly. And the $x$-mean error,
$\Delta:=\bar X-\hat X$, satisfies

\[
\Delta=\frac1N\sum_i c_i\,\delta_i,
\qquad
|\Delta|\le \frac1N\sum_i c_i\,\frac{w_i}{2}\;=:\;D.
\]
\end{lemma}

\begin{proof}
Group the global $y$-average by bucket. The points in bucket $i$
contribute total $y$ equal to $c_i\bar y_i=c_i\mu_i$, so

\[
\bar Y=\frac1N\sum_{i,j}y_{ij}
      =\frac1N\sum_i c_i\bar y_i
      =\frac1N\sum_i c_i\mu_i
      =\hat Y.
\]

This is exact -- nothing was approximated, because $\mu_i$ is the exact
bucket mean. The $y$ side carries no slack.

For $x$, do the same grouping but substitute
$\bar x_i = m_i + \delta_i$ (the definition of $\delta_i$):

\[
\bar X=\frac1N\sum_i c_i\bar x_i
      =\frac1N\sum_i c_i\,(m_i+\delta_i)
      =\underbrace{\frac1N\sum_i c_i m_i}_{=\,\hat X}
      \;+\;\frac1N\sum_i c_i\delta_i .
\]

So $\bar X-\hat X=\frac1N\sum_i c_i\delta_i$, which is the stated
$\Delta$. Bounding it: by the triangle inequality and $(\star)$,

\[
|\Delta|\le\frac1N\sum_i c_i\,|\delta_i|
        \le\frac1N\sum_i c_i\,\frac{w_i}{2}=D. \qedhere
\]
\end{proof}

\noindent Plain meaning: our $y$ stand-in is perfect; our $x$ stand-in
is off by an unknown amount $\Delta$, but that amount is provably no
larger than $D$, a number we can compute from widths and counts alone.

\section{Step 2 --- separating ``between'' and ``within''}

To bound $S_{xx}$ and $S_{xy}$ we split each into two parts: variation
\emph{between} the buckets (driven by where the bucket centroids sit)
and variation \emph{within} each bucket (driven by the spread of points
around their own centroid). This split is exact.

For any bucket the points satisfy two zero-sum identities by definition
of the centroid:

\[
\sum_{j}(x_{ij}-\bar x_i)=0,
\qquad
\sum_{j}(y_{ij}-\bar y_i)=0.
\]

Write $x_{ij}-\bar X=(\bar x_i-\bar X)+(x_{ij}-\bar x_i)$ and similarly
for $y$, then expand the squares/products in $S_{xx},S_{xy}$. Every
cross term contains a factor $\sum_j(x_{ij}-\bar x_i)$ or
$\sum_j(y_{ij}-\bar y_i)$, which is zero, so the cross terms vanish and
we are left with exactly:

\[
S_{xx}=\underbrace{\sum_i c_i(\bar x_i-\bar X)^2}_{\text{between}}
       \;+\;\underbrace{\sum_i A_i}_{\text{within}},
\qquad
A_i=\sum_j (x_{ij}-\bar x_i)^2,
\]

\[
S_{xy}=\underbrace{\sum_i c_i(\bar x_i-\bar X)(\mu_i-\bar Y)}_{\text{between}}
       \;+\;\underbrace{\sum_i B_i}_{\text{within}},
\qquad
B_i=\sum_j (x_{ij}-\bar x_i)(y_{ij}-\bar y_i).
\]

Both identities are exact. Now we bound the two parts separately: the
between part using Lemma~\ref{lem:means} (it involves the unknown
$\bar X$, which we replace by $\hat X$ at the cost of $\Delta$), and the
within part using the bucket's min/max and width.

\section{Step 3 --- bounding $S_{xy}$ (Lemma B)}

Define the quantity we \emph{can} compute -- the between part with the
unknown true means replaced by the computable stand-ins:

\[
P_0=\sum_i c_i\,(m_i-\hat X)(\mu_i-\hat Y).
\]

\begin{lemma}\label{lem:B}
$\bigl|\,S_{xy}-P_0\,\bigr|\le E_P$, where, writing
$d_i=\max(h_i-\mu_i,\ \mu_i-\ell_i)$ for the largest possible
$y$-deviation in bucket $i$,

\[
E_P=\sum_i c_i\,|\mu_i-\hat Y|\,\frac{w_i}{2}
   \;+\;\sum_i c_i\,w_i\,d_i.
\]
\end{lemma}

\begin{proof}
Start from the exact between part of $S_{xy}$ and replace the unknown
$\bar x_i-\bar X$ by computable pieces plus error. Using
$\bar x_i=m_i+\delta_i$ and $\bar X=\hat X+\Delta$ (Lemma~A),

\[
\bar x_i-\bar X=(m_i-\hat X)+(\delta_i-\Delta).
\]

Also $\bar Y=\hat Y$ exactly (Lemma~A), so $\mu_i-\bar Y=\mu_i-\hat Y$.
Substitute into the between sum:

\[
\sum_i c_i(\bar x_i-\bar X)(\mu_i-\bar Y)
=\sum_i c_i\bigl[(m_i-\hat X)+(\delta_i-\Delta)\bigr](\mu_i-\hat Y).
\]

Distribute into three groups:

\[
=\underbrace{\sum_i c_i(m_i-\hat X)(\mu_i-\hat Y)}_{=\,P_0}
\;+\;\sum_i c_i\,\delta_i(\mu_i-\hat Y)
\;-\;\Delta\sum_i c_i(\mu_i-\hat Y).
\]

The last group vanishes. Why: $\Delta$ does not depend on $i$, so pull
it out, and

\[
\sum_i c_i(\mu_i-\hat Y)
=\Bigl(\sum_i c_i\mu_i\Bigr)-\hat Y\sum_i c_i
=N\hat Y-\hat Y\,N
=0,
\]

using the definition $\hat Y=\frac1N\sum_i c_i\mu_i$. So the entire
$\Delta$ term is $-\Delta\cdot 0=0$. This is the key cancellation:
the error in our $x$-mean stand-in does \emph{not} corrupt the between
part of $S_{xy}$ at first order.

What remains, after adding back the within part $\sum_i B_i$:

\[
S_{xy}-P_0
=\underbrace{\sum_i c_i\,\delta_i(\mu_i-\hat Y)}_{\text{(I): unknown }x\text{-placement}}
\;+\;\underbrace{\sum_i B_i}_{\text{(II): within-bucket spread}}.
\]

Bound each piece using only what the summaries give us.

\emph{Piece (I).} The only unknown is $\delta_i$, and $(\star)$ says
$|\delta_i|\le w_i/2$. Everything else is computable. So

\[
\Bigl|\sum_i c_i\,\delta_i(\mu_i-\hat Y)\Bigr|
\le\sum_i c_i\,|\mu_i-\hat Y|\,|\delta_i|
\le\sum_i c_i\,|\mu_i-\hat Y|\,\frac{w_i}{2}.
\]

That is exactly the first sum in $E_P$.

\emph{Piece (II).} Inside bucket $i$, bound the product term by term.
Both $x_{ij}$ and $\bar x_i$ lie in the span $[a_i,b_i]$, so their
difference is at most the width: $|x_{ij}-\bar x_i|\le w_i$. And every
$y_{ij}$ lies in $[\ell_i,h_i]$ while $\bar y_i=\mu_i$, so
$|y_{ij}-\mu_i|\le\max(h_i-\mu_i,\mu_i-\ell_i)=d_i$. Therefore

\[
|B_i|\le\sum_j |x_{ij}-\bar x_i|\,|y_{ij}-\mu_i|
\le\sum_j w_i\,d_i
= c_i\,w_i\,d_i,
\]

and summing over buckets gives the second sum in $E_P$.

Adding the two bounds via the triangle inequality yields
$|S_{xy}-P_0|\le E_P$.
\end{proof}

\noindent Plain meaning: $P_0$ is our computable estimate of $S_{xy}$.
It is wrong by at most $E_P$, and $E_P$ has two clearly identified
causes -- not knowing the horizontal placement (piece I) and not
knowing how points scatter inside a bucket (piece II). Both are
controlled purely by widths, counts, and the stored min/max.

\section{Step 4 --- bounding $S_{xx}$ (Lemma C)}

Now the denominator. Define the computable estimate

\[
Q_0=\sum_i c_i\,u_i^2,
\qquad u_i:=m_i-\hat X .
\]

\begin{lemma}\label{lem:C}
$Q_0-E_\downarrow\;\le\;S_{xx}\;\le\;Q_0+E_\downarrow+E_\uparrow$, where

\[
E_\downarrow=\sum_i c_i\,|m_i-\hat X|\,w_i,
\qquad
E_\uparrow=\sum_i c_i\Bigl(\tfrac{w_i}{2}+D\Bigr)^2
          +\sum_i c_i\,\tfrac{w_i^2}{4}.
\]
\end{lemma}

\begin{proof}
Take the exact between part of $S_{xx}$ and substitute
$\bar x_i-\bar X=u_i+e_i$ where $e_i:=\delta_i-\Delta$ collects all the
$x$-placement uncertainty. Expanding the square:

\[
\sum_i c_i(\bar x_i-\bar X)^2
=\sum_i c_i u_i^2
\;+\;2\sum_i c_i u_i e_i
\;+\;\sum_i c_i e_i^2 .
\]

The first sum is $Q_0$. In the cross sum, replace $e_i$ by
$\delta_i-\Delta$; the $-\Delta$ piece drops out because
$\sum_i c_i u_i=\sum_i c_i(m_i-\hat X)=N\hat X-\hat X N=0$ (same kind of
cancellation as before). So $\sum_i c_i u_i e_i=\sum_i c_i u_i\delta_i$.
Putting the within part $\sum_i A_i$ back:

\[
S_{xx}=Q_0
\;+\;2\sum_i c_i u_i\delta_i
\;+\;\underbrace{\sum_i c_i e_i^2}_{\ge 0}
\;+\;\underbrace{\sum_i A_i}_{\ge 0}.
\tag{$\dagger$}
\]

The last two sums are sums of squares, hence $\ge 0$. The middle
(signed) term is controlled by $(\star)$:

\[
\Bigl|2\sum_i c_i u_i\delta_i\Bigr|
\le 2\sum_i c_i\,|u_i|\,|\delta_i|
\le 2\sum_i c_i\,|m_i-\hat X|\,\frac{w_i}{2}
=\sum_i c_i\,|m_i-\hat X|\,w_i
=E_\downarrow.
\]

\emph{Lower bound.} In $(\dagger)$, the two squared sums only add to
$S_{xx}$, so dropping them can only decrease the right-hand side:

\[
S_{xx}\ge Q_0+2\sum_i c_i u_i\delta_i\ge Q_0-E_\downarrow.
\]

\emph{Upper bound.} Bound the two non-negative sums from above.

For $\sum_i c_i e_i^2$: $e_i=\delta_i-\Delta$, and
$|\delta_i|\le w_i/2$ by $(\star)$ while $|\Delta|\le D$ by Lemma~A, so
$|e_i|\le w_i/2+D$ and $e_i^2\le(w_i/2+D)^2$. Hence
$\sum_i c_i e_i^2\le\sum_i c_i(w_i/2+D)^2$, the first sum in
$E_\uparrow$.

For $\sum_i A_i$: $A_i=\sum_j(x_{ij}-\bar x_i)^2$ is the total squared
spread of $c_i$ points confined to an interval of width $w_i$. That
spread is largest when the points are pushed as far apart as the
interval allows, i.e.\ split between the two endpoints; a short
calculation shows the maximum is $c_i w_i^2/4$. So
$A_i\le c_i w_i^2/4$ and $\sum_i A_i\le\sum_i c_i w_i^2/4$, the second
sum in $E_\uparrow$.

Substituting both into $(\dagger)$ and using
$|2\sum c_i u_i\delta_i|\le E_\downarrow$:

\[
S_{xx}\le Q_0+E_\downarrow+E_\uparrow. \qedhere
\]
\end{proof}

\noindent Plain meaning: $Q_0$ is our computable estimate of $S_{xx}$.
Unlike $S_{xy}$, the uncertainty here is one-sided in character -- the
unknown placement and the within-bucket spread can only \emph{inflate}
$S_{xx}$ relative to the estimate (that is why the upper error
$E_\uparrow$ is separate and larger), while the lower side is
controlled by the single signed term $E_\downarrow$.

\section{Step 5 --- from the two boxes to a slope interval}

We now know, with certainty,

\[
S_{xy}\in[\,P^-,P^+\,],\quad P^\pm=P_0\pm E_P,
\qquad
S_{xx}\in[\,Q^-,Q^+\,],\quad
\begin{aligned}
Q^-&=Q_0-E_\downarrow,\\
Q^+&=Q_0+E_\downarrow+E_\uparrow.
\end{aligned}
\]

So the true pair $(S_{xy},S_{xx})$ lies somewhere in the rectangle
$[P^-,P^+]\times[Q^-,Q^+]$. We want the range of the ratio
$\beta=S_{xy}/S_{xx}$ over that rectangle.

\begin{theorem}
Suppose $Q^->0$ (the routine checks this; otherwise see the remark
below). Let

\[
C=\Bigl\{\,\tfrac{P^-}{Q^-},\ \tfrac{P^-}{Q^+},\
            \tfrac{P^+}{Q^-},\ \tfrac{P^+}{Q^+}\,\Bigr\}
\]

be the four ratios at the rectangle's corners. Then

\[
\min C\;\le\;\beta\;\le\;\max C.
\]
\end{theorem}

\begin{proof}
Consider $f(p,q)=p/q$ on the rectangle. Because $Q^->0$, the whole
rectangle has $q>0$, so $f$ is finite throughout. Hold $q$ fixed:
$f$ is linear in $p$, hence monotone in $p$. Hold $p$ fixed:
$\partial f/\partial q=-p/q^2$ has a constant sign on the rectangle
(the sign of $-p$, and $p$ does not change sign because $[P^-,P^+]$ is
one interval), hence $f$ is monotone in $q$. A function that is
monotone in each coordinate separately on a rectangle attains its
maximum and minimum at the rectangle's corners. So the extreme values
of $f$ over the rectangle are among the four corner ratios in $C$.
Since the true $(S_{xy},S_{xx})$ is some point of the rectangle, its
ratio $\beta$ is between the smallest and largest of those corner
ratios.
\end{proof}

\paragraph{The degenerate case.}
If $Q^-\le 0$ (or the estimate $Q_0\le 0$), the denominator interval
$[Q^-,Q^+]$ straddles $0$. Then $\beta$ is genuinely unbounded -- some
point clouds consistent with the summaries make $S_{xx}$ arbitrarily
small and the slope blows up -- so the only honest answer is
$(-\infty,+\infty)$, and that is what the routine returns. It is a
valid (if useless) enclosure.

\section{What the bound is, in one paragraph}

We could not compute the slope because the raw points were discarded
and many point clouds fit the same summaries. We replaced the two
unknown means by computable midpoint surrogates (Lemma~A), split each
of $S_{xx},S_{xy}$ into an exact between/within decomposition, and
bounded every place an unknown entered using only two facts: each
point's $x$ sits within its bucket's width $w_i$ of the centroid
($\star$), and each point's $y$ sits within the stored
$[\ell_i,h_i]$. That produced a computable estimate and a computable
error bar for the numerator ($P_0\pm E_P$) and for the denominator
($Q_0$, with $-E_\downarrow$ / $+E_\downarrow+E_\uparrow$). The true
$(S_{xy},S_{xx})$ lives in the resulting rectangle, and a ratio is
monotone on a positive rectangle, so the four corner ratios bracket
the true slope. Every error term is proportional to a bucket
\emph{width}: shrink the buckets and the interval shrinks with them.
That is the whole proof, and the practical lever.

\paragraph{Numerical note.}
Everything above is exact in real arithmetic. A floating-point
implementation should widen the final interval by the accumulated
round-off (roughly machine epsilon times the magnitude of the sums).
That is slack, not a gap in the argument.

\end{document}
